\input{../tex-inputs/latex-korrekturansicht-vorspann}

\section[Theodor Herzl an Arthur Schnitzler, 4. 5. 1893]{L03827 Theodor Herzl an Arthur Schnitzler, 4. 5. 1893}
\nopagebreak\mylabel{L03827v}
\rehead{ }\normalsize\beginnumbering\briefempfaengerindex{Schnitzler, Arthur@\textsc{Schnitzler, Arthur}!zzzHerzl, Theodor@\emph{von Theodor Herzl}!1893-05-041@{4. 5. 1893}|(be}
\toendnotes[C]{\smallbreak\pagebreak[2]}
\correspDesc{Versand  durch Theodor Herzl am 4. 5. 1893 in Paris
\newline{}Erhalt  durch Arthur Schnitzler im Zeitraum [5. 5. 1893
                  – 9. 5. 1893?] in Wien}\toendnotes[C]{\smallbreak}
\Standort{CUL, Schnitzler, B 39.}
\physDesc{Brief, 1 Blatt, 1 Seite, 535 Zeichen
\newline{}Handschrift: schwarze Tinte, lateinische Kurrent
\newline{}Ordnung: mit Bleistift von unbekannter Hand nummeriert: »8« }
\buchAbdrucke{\weitereDrucke{Theodor Herzl: \emph{Briefe und
                        autobiographische Notizen 1866–1895}. Bearbeitet von Johannes Wachten in Zusammenarbeit mit Chaya Harel, Daisy Tycho und Manfred Winkler. Berlin, Frankfurt am Main, Wien: \emph{Propyläen} 1983, S. 526 (Briefe und Tagebücher. Herausgegeben von Alex Bein, Hermann Greive, Moshe Schaerf, Julius H. Schoeps und Johannes Wachten, 1).} }\toendnotes[C]{\smallbreak}
\pstart
           {\pb}\textcolor{gray}{\textbf{\textsc{\textcolor{brown}{Nouvelle Presse Libre}\orgindex{Neue Freie Presse@Neue Freie Presse|pw}{}\ledrightnote{\textcolor{brown}{Neue Freie Presse}}}}}\hfill \textcolor{pink}{\textcolor{gray}{\textbf{8, RUE DE MONCEAU}}}\oindex{8, rue de Monceau@\textbf{8, rue de Monceau}, \emph{Wohngebäude}|pw}{}\ledrightnote{\textcolor{pink}{8, rue de Monceau}}\pend
           
\pstart
           \textcolor{gray}{\textbf{D\textsuperscript{r}{ }TH. HERZL}}\pend
           
\pstart{}Mein lieber Freund!\pend\vspace{0.5em}
\pstart
           Sehr erschüttert lese ich \label{K_L03827-1v}\edtext{in der
                  Zeitung}{\lemma{\textnormal{\emph{in der
                  Zeitung}}}\Cendnote{\textnormal{Johann Schnitzler starb am
                     2. 5. 1893. Am Folgetag wurde eine Traueranzeige in der \emph{\textcolor{green}{Presse}\pwindex{Presse@\emph{Die Presse}|pwk}}, der \emph{\textcolor{green}{Neuen Freien Presse}\pwindex{Neue Freie Presse@\emph{Neue Freie Presse}|pwk}} und der \emph{\textcolor{green}{Wiener Zeitung}\pwindex{Wiener Zeitung@\emph{Wiener Zeitung}|pwk}} gedruckt. Einen Nachruf brachte etwa die \emph{\textcolor{green}{Wiener Allgemeine Zeitung}\pwindex{Wiener Allgemeine Zeitung@\emph{Wiener Allgemeine Zeitung}|pwk}} (\emph{\textcolor{green}{Professor Dr. Johann Schnitzler}\pwindex{Professor Dr. Johann Schnitzler@\emph{Professor Dr. Johann Schnitzler}|pwk}}. In: \emph{\textcolor{green}{Wiener Allgemeine Zeitung}\pwindex{Wiener Allgemeine Zeitung@\emph{Wiener Allgemeine Zeitung}|pwk}}, Nr. 4528,
                        3. 5. 1893, S. 6).}}}\label{K_L03827-1}, dass Ihr \textcolor{blue}{Vater}\pwindex{Schnitzler, Johann 10.\,4.\,1835 Nagykanizsa – 2.\,5.\,1893 Wien@\textsc{Schnitzler, Johann} (10.\,4.\,1835 Nagykanizsa – 2.\,5.\,1893 Wien), \emph{Laryngologe}|pwv}{}\ledrightnote{{$\rightarrow$}\emph{\textcolor{blue}{Johann Schnitzler}}} gestorben ist.\pend
           
\pstart
           Wie arm ist unsere Rede, wenn wir einen wirklichen grossen Schmerz vor uns haben.\pend
           
\pstart
           Ein stummer Händedruck sagt die Theilnahme noch am besten – lassen Sie diese Zeilen
               dafür gelten.\pend
           
\pstart
           Sagen Sie auch Ihrer verehrten Frau \textcolor{blue}{Mutter}\pwindex{Schnitzler, Louise 8.\,7.\,1840 Kőszeg – 9.\,9.\,1911 Wien@\textsc{Schnitzler, Louise} (8.\,7.\,1840 Kőszeg – 9.\,9.\,1911 Wien)|pwv}{}\ledrightnote{{$\rightarrow$}\emph{\textcolor{blue}{Louise Schnitzler}}} und Ihren lieben \textcolor{blue}{Geschwistern}\pwindex{Schnitzler, Julius 13.\,7.\,1865 Wien – 29.\,6.\,1939 ebd.@\textsc{Schnitzler, Julius} (13.\,7.\,1865 Wien – 29.\,6.\,1939 ebd.), \emph{Chirurg}|pwv}\pwindex{Hajek, Gisela 20.\,12.\,1867 Wien – 3.\,2.\,1953 Cambridge@\textsc{Hajek, Gisela} (20.\,12.\,1867 Wien – 3.\,2.\,1953 Cambridge)|pwv}{}\ledrightnote{{$\rightarrow$}\emph{\textcolor{blue}{Julius Schnitzler}}{\newline}{$\rightarrow$}\emph{\textcolor{blue}{Gisela Hajek}}}, dass ich zu denen
               gehöre, die an Ihrem schweren Verlust \introOben{}am\introOben{} Innigsten
               theilnehmen.\pend
           
\pstart
           Leben Sie wohl, mein lieber Schnitzler und glauben Sie an die Freundschaft
               {\\[\baselineskip]}Ihres herzlich ergebenen {\\[\baselineskip]}\spacefill\mbox{Th Herzl}\pend
           \leftskip=0em{}
\pstart
           \textcolor{pink}{Paris}\oindex{Paris@\textbf{Paris}, \emph{Hauptstadt}|pw}{}\ledrightnote{\textcolor{pink}{Paris}}{ }4 Mai 93\pend
           \selectlanguage{ngerman}\endnumbering\briefempfaengerindex{Schnitzler, Arthur@\textsc{Schnitzler, Arthur}!zzzHerzl, Theodor@\emph{von Theodor Herzl}!1893-05-041@{4. 5. 1893}|)be}\mylabel{L03827h}
\begin{anhang}
\end{anhang}\newcommand{\dateiname}{L03827}\newcommand{\titel}{Theodor Herzl an Arthur Schnitzler, 4. 5. 1893}\newcommand{\editorInnen}{Selma Jahnke und Martin Anton Müller}\input{../tex-inputs/latex-korrekturansicht-abspann}
